\chapter{Resultados e discussão}
\label{cap:resultados}

\section{Exemplificação do uso do artefacto}
\label{sec:caso}

\subsection{Caso CONCEPTIO --- loop fechado SysON $\leftrightarrow$ MQTT}

Demonstração ponta a ponta: modelo com múltiplos Success Criteria (threshold,
range, statistical, \textbf{rate}, \textbf{dwell}) no SysON; import e gate no
ReqValLive; medição com telemetria MQTT do laboratório ou simulador
(\texttt{scripts/publish\_three\_drones.py}); publicação de
\texttt{VerificationResult} no modelo.

O orientador validou o ciclo duro em agosto/2026 (SysON + MQTT + write-back).
Este capítulo documenta o ensaio formal com critérios Jefferson de taxa e
persistência, descrito em \texttt{docs/ENSAIO\_SYSON\_RATE\_DWELL.md}.

\begin{table}[htbp]
  \centering
  \caption{Requisitos do demo SysON usados no ensaio rate/dwell (2026-08).}
  \label{tab:ensaio-reqs}
  \begin{tabular}{@{}lllp{4.2cm}@{}}
    \toprule
    ID & Tipo SC & Veredicto esperado & Nota \\
    \midrule
    \texttt{RQ\_BAT\_RATE\_001} & rate & FAIL & $\approx$60\%/min a 1\,Hz no simulador \\
    \texttt{RQ\_BAT\_DWELL\_001} & dwell & FAIL & charlie inicia a 15\% ($<$20\%) \\
    \texttt{RQ\_BAT\_001} & threshold & PASS/FAIL & conforme corrida \\
    \texttt{RQ\_SEP\_001} & threshold global & PASS & separação Haversine \\
    \bottomrule
  \end{tabular}
\end{table}

% TODO: inserir figuras (explorador SysON, UI live, laudo HTML) — ver ENSAIO doc.

\subsection{Protocolo de ensaio}

\begin{enumerate}
  \item Importar \texttt{reqvallive\_demo.sysml} no SysON; Documentation =
        só \texttt{\_go\_to\_verification}; SC no item \texttt{SuccessCriteria}.
  \item ReqValLive: import $\rightarrow$ gate ACCEPT $\rightarrow$ montar GSE.
  \item Ligar MQTT (broker/tópico do lab ou simulador local).
  \item Iniciar medição ($\geq$30\,s); encerrar; gerar laudo HTML (secções de
        procedimento: método Test, tópico, janela, snapshot SC).
  \item Publicar \texttt{VerificationResult} no SysON; capturas para este capítulo.
\end{enumerate}

\section{Discussão do uso do artefacto}
\label{sec:discussao}

\subsection{O que o ciclo resolve}

Fecha o gap \emph{declarar SC $\rightarrow$ medir $\rightarrow$ laudar $\rightarrow$
escrever no modelo} sem depender do runtime de simulação SysML no Magic.

\subsection{Naturezas de observação --- cobertura actual vs.\ roadmap}

A taxonomia GSE (Jefferson) distingue valor, tempo, sequência, taxa, eventos,
frequência e consistência entre estados~\cite{jefferson2026gse}.
A PoC implementa \textbf{quatro} naturezas no motor live:
\begin{itemize}
  \item \textbf{valor} --- threshold e range;
  \item \textbf{amplitude em janela} --- statistical (range/max/min);
  \item \textbf{taxa} --- \texttt{rate} ($\Delta/\Delta t$);
  \item \textbf{persistência} --- \texttt{dwell} (condição por tempo máximo).
\end{itemize}
Sequência de eventos, ocorrência/ausência, frequência e estados permanecem
trabalho futuro --- telemetria do lab ainda não as fornece de forma estruturada.

\subsection{Ameaças à validade}

\begin{itemize}
  \item Escala de laboratório (simulador + voos reais pontuais).
  \item Dependência do protocolo MQTT do CONCEPTIO (JSON por drone).
  \item Expectativa errada sobre SysML (esqueleto vs.\ simulação) --- mitigada
        na UI e no laudo.
  \item Taxa derivada sensível ao intervalo de amostragem MQTT.
\end{itemize}

\section{Verificação dos objetivos}
\label{sec:objetivos-fechados}

\begin{table}[htbp]
  \centering
  \caption{Objectivos específicos vs.\ evidência (2026-08).}
  \label{tab:objetivos}
  \begin{tabular}{@{}lp{7.2cm}@{}}
    \toprule
    Objectivo & Estado / evidência \\
    \midrule
    Gate ACCEPT/REJECT & Implementado; demo orientador ago/2026 \\
    Avaliação live MQTT & threshold, range, statistical, \textbf{rate}, \textbf{dwell}; latch FAIL \\
    Laudo + esqueleto SysML & HTML com procedimento + snapshot SC; VerificationResult \\
    UI principal em inglês & SysON como fluxo principal; CATIA legado secundário \\
    Ensaio documentado rate/dwell & \texttt{docs/ENSAIO\_SYSON\_RATE\_DWELL.md} \\
    \bottomrule
  \end{tabular}
\end{table}

Regressão: 113 testes \texttt{pytest} (incl.\ \texttt{test\_jefferson\_natures.py}).
